% GENERATED FILE. Edit canonical lessons, then run npm run generate:books.
% canonical-source-hash: fnv1a64:b9b2e8453ae5d6d0
% canonical-lessons: GU-C386-astha, GU-C386-asro, GU-C386-agrah, GU-C386-adhin, GU-C386-jhukvu

\chapter{Trust and Dependence}
\label{ch:gu-trust-and-dependence}
\hlchaptermodality{386}

\begin{chapteropening}
\textbf{By the end of this chapter:} \emph{I can say what will happen if someone leans on others.}

Complete the last lesson of chapter 386: I can say what will happen if someone leans on others.
\end{chapteropening}

\section[\texorpdfstring{\gu{આસ્થા} (āsthā)}{આસ્થા (āsthā)}]{\gu{આસ્થા} (āsthā) — belief; faith (in God or an idea)}

\label{lesson:GU-C386-astha}



\begin{quote}
Before the new one: say the Gujarati for to attain, then the Gujarati for to honour with an award.
\end{quote}

\subsection*{You'll want to know: \gu{આસ્થા}}
\textbf{\gu{આસ્થા}} — \emph{āsthā} — \textquotedblleft{}belief; faith (in God or an idea)\textquotedblright{}.

\begin{quote}
\gu{જો} \gu{તમારી} \gu{આસ્થા} \gu{મજબૂત} \gu{હશે} \gu{તો} \gu{તમે} \gu{હિંમત} \gu{નહીં} \gu{હારો}.
\end{quote}

\begin{quote}
\emph{jo tamārī āsthā majbūt haśe to tame himmat nahī̃ hāro.}
\end{quote}

\begin{quote}
If your faith is strong, you will not lose heart.
\end{quote}

A feminine noun.

\begin{cousinweb}
\textbf{\gu{આસ્થા}} (\emph{āsthā}) is a tatsama, borrowed from Sanskrit \emph{āsthā}, \textquotedblleft{}regard, care, confidence\textquotedblright{} (Monier-Williams).
\end{cousinweb}

\begin{grammarlens}[title={What you've built}]
One more word for saying what happens if.
\end{grammarlens}

\subsection*{Your turn}
\begin{itemize}
\raggedright
  \item \emph{Say it:} \emph{āsthā}
  \item \emph{Say it:} \emph{āsthā}, once more
  \item \emph{Say it:} \emph{jo tamārī āsthā majbūt haśe to tame himmat nahī̃ hāro.} — If your faith is strong, you will not lose heart.
  \item \emph{Recall:} say the Gujarati for to attain, then the Gujarati for to honour with an award, then say \emph{āsthā} again
\end{itemize}

\begin{tcolorbox}[breakable,colback=teal!4,colframe=teal!35!black,title={Before you move on}]
What does \gu{આસ્થા} mean? (\textquotedblleft{}Belief; faith (in God or an idea)\textquotedblright{}.) Say it once more.
\end{tcolorbox}
\section[\texorpdfstring{\gu{આશરો} (āśro)}{આશરો (āśro)}]{\gu{આશરો} (āśro) — a refuge}

\label{lesson:GU-C386-asro}



\begin{quote}
Before the new one: say the Gujarati for to honour with an award, then the Gujarati for belief.
\end{quote}

\subsection*{You'll want to know: \gu{આશરો}}
\textbf{\gu{આશરો}} — \emph{āśro} — \textquotedblleft{}a refuge\textquotedblright{}.

\begin{quote}
\gu{જો} \gu{પૂર} \gu{આવશે} \gu{તો} \gu{આખું} \gu{ગામ} \gu{શાળામાં} \gu{આશરો} \gu{લેશે}.
\end{quote}

\begin{quote}
\emph{jo pūr āvśe to ākhũ gām śāḷāmā̃ āśro leśe.}
\end{quote}

\begin{quote}
If the flood comes, the whole village will take shelter in the school.
\end{quote}

A masculine noun.

\begin{cousinweb}
\textbf{\gu{આશરો}} (\emph{āśro}) CDIAL lists the Gujarati noun, \textquotedblleft{}protection\textquotedblright{}, under a reconstructed \emph{āśara} (CDIAL 1455).
\end{cousinweb}

\begin{grammarlens}[title={What you've built}]
One more word for saying what happens if.
\end{grammarlens}

\subsection*{Your turn}
\begin{itemize}
\raggedright
  \item \emph{Say it:} \emph{āśro}
  \item \emph{Say it:} \emph{āśro}, once more
  \item \emph{Say it:} \emph{jo pūr āvśe to ākhũ gām śāḷāmā̃ āśro leśe.} — If the flood comes, the whole village will take shelter in the school.
  \item \emph{Recall:} say the Gujarati for to honour with an award, then the Gujarati for belief, then say \emph{āśro} again
\end{itemize}

\begin{tcolorbox}[breakable,colback=teal!4,colframe=teal!35!black,title={Before you move on}]
What does \gu{આશરો} mean? (\textquotedblleft{}A refuge\textquotedblright{}.) Say it once more.
\end{tcolorbox}
\section[\texorpdfstring{\gu{આગ્રહ} (āgrah)}{આગ્રહ (āgrah)}]{\gu{આગ્રહ} (āgrah) — insistence}

\label{lesson:GU-C386-agrah}



\begin{quote}
Before the new one: say the Gujarati for belief, then the Gujarati for a refuge.
\end{quote}

\subsection*{You'll want to know: \gu{આગ્રહ}}
\textbf{\gu{આગ્રહ}} — \emph{āgrah} — \textquotedblleft{}insistence\textquotedblright{}.

\begin{quote}
\gu{જો} \gu{તમે} \gu{આગ્રહ} \gu{કરશો} \gu{તો} \gu{હું} \gu{જમીને} \gu{જઈશ}.
\end{quote}

\begin{quote}
\emph{jo tame āgrah karśo to hũ jamīne jaīś.}
\end{quote}

\begin{quote}
If you insist, I will eat before I go.
\end{quote}

A masculine noun.

\begin{cousinweb}
\textbf{\gu{આગ્રહ}} (\emph{āgrah}) is a tatsama, borrowed from Sanskrit \emph{āgraha}, \textquotedblleft{}insisting on, a strong inclination\textquotedblright{} (Monier-Williams).
\end{cousinweb}

\begin{grammarlens}[title={What you've built}]
One more word for saying what happens if.
\end{grammarlens}

\subsection*{Your turn}
\begin{itemize}
\raggedright
  \item \emph{Say it:} \emph{āgrah}
  \item \emph{Say it:} \emph{āgrah}, once more
  \item \emph{Say it:} \emph{jo tame āgrah karśo to hũ jamīne jaīś.} — If you insist, I will eat before I go.
  \item \emph{Recall:} say the Gujarati for belief, then the Gujarati for a refuge, then say \emph{āgrah} again
\end{itemize}

\begin{tcolorbox}[breakable,colback=teal!4,colframe=teal!35!black,title={Before you move on}]
What does \gu{આગ્રહ} mean? (\textquotedblleft{}Insistence\textquotedblright{}.) Say it once more.
\end{tcolorbox}
\section[\texorpdfstring{\gu{આધીન} (ādhīn)}{આધીન (ādhīn)}]{\gu{આધીન} (ādhīn) — subject to; dependent on}

\label{lesson:GU-C386-adhin}



\begin{quote}
Before the new one: say the Gujarati for a refuge, then the Gujarati for insistence.
\end{quote}

\subsection*{You'll want to know: \gu{આધીન}}
\textbf{\gu{આધીન}} — \emph{ādhīn} — \textquotedblleft{}subject to; dependent on\textquotedblright{}.

\begin{quote}
\gu{આ} \gu{યોજના} \gu{સરકારની} \gu{મંજૂરીને} \gu{આધીન} \gu{છે}; \gu{જો} \gu{મંજૂરી} \gu{આવશે} \gu{તો} \gu{કામ} \gu{શરૂ} \gu{થશે}.
\end{quote}

\begin{quote}
\emph{ā yojnā sarkārnī manjūrīne ādhīn che; jo manjūrī āvśe to kām śarū thaśe.}
\end{quote}

\begin{quote}
This plan is subject to government approval; if approval comes, work will start.
\end{quote}

\begin{cousinweb}
\textbf{\gu{આધીન}} (\emph{ādhīn}) is a tatsama, borrowed from Sanskrit \emph{ādhīna}, \textquotedblleft{}dependent on, subject to\textquotedblright{} (Monier-Williams).
\end{cousinweb}

\begin{grammarlens}[title={What you've built}]
One more word for saying what happens if.
\end{grammarlens}

\subsection*{Your turn}
\begin{itemize}
\raggedright
  \item \emph{Say it:} \emph{ādhīn}
  \item \emph{Say it:} \emph{ādhīn}, once more
  \item \emph{Say it:} \emph{ā yojnā sarkārnī manjūrīne ādhīn che; jo manjūrī āvśe to kām śarū thaśe.} — This plan is subject to government approval; if approval comes, work will start.
  \item \emph{Recall:} say the Gujarati for a refuge, then the Gujarati for insistence, then say \emph{ādhīn} again
\end{itemize}

\begin{tcolorbox}[breakable,colback=teal!4,colframe=teal!35!black,title={Before you move on}]
What does \gu{આધીન} mean? (\textquotedblleft{}Subject to; dependent on\textquotedblright{}.) Say it once more.
\end{tcolorbox}
\section[\texorpdfstring{\gu{ઝૂકવું} (jhūkvũ)}{ઝૂકવું (jhūkvũ)}]{\gu{ઝૂકવું} (jhūkvũ) — to give in; to bend}

\label{lesson:GU-C386-jhukvu}



\begin{quote}
Before the new one: say the Gujarati for insistence, then the Gujarati for subject to.
\end{quote}

\subsection*{You'll want to know: \gu{ઝૂકવું}}
\textbf{\gu{ઝૂકવું}} — \emph{jhūkvũ} — \textquotedblleft{}to give in; to bend\textquotedblright{}.

\begin{quote}
\gu{જો} \gu{આપણે} \gu{બધા} \gu{સાથે} \gu{રહીશું} \gu{તો} \gu{સરકારે} \gu{ઝૂકવું} \gu{પડશે}.
\end{quote}

\begin{quote}
\emph{jo āpaṇe badhā sāthe rahīśũ to sarkāre jhūkvũ paḍśe.}
\end{quote}

\begin{quote}
If we all stand together, the government will have to give in.
\end{quote}

Past (he …): \gu{તે} \gu{છેવટે} \gu{ઝૂક્યો}, \emph{te chevṭe jhūkyo}.

\begin{cousinweb}
\textbf{\gu{ઝૂકવું}} (\emph{jhūkvũ}) CDIAL lists the Gujarati verb, \textquotedblleft{}to bend\textquotedblright{}, under a reconstructed \emph{jhukkati}, \textquotedblleft{}stoops\textquotedblright{} (CDIAL 5399).
\end{cousinweb}

\begin{grammarlens}[title={What you've built}]
One more word for saying what happens if.
\end{grammarlens}

\subsection*{Your turn}
\begin{itemize}
\raggedright
  \item \emph{Say it:} \emph{jhūkvũ}
  \item \emph{Say it:} \emph{jhūkvũ}, once more
  \item \emph{Say it:} \emph{jo āpaṇe badhā sāthe rahīśũ to sarkāre jhūkvũ paḍśe.} — If we all stand together, the government will have to give in.
  \item \emph{Recall:} say the Gujarati for insistence, then the Gujarati for subject to, then say \emph{jhūkvũ} again
\end{itemize}

\begin{tcolorbox}[breakable,colback=teal!4,colframe=teal!35!black,title={Before you move on}]
What does \gu{ઝૂકવું} mean? (\textquotedblleft{}To give in; to bend\textquotedblright{}.) Say it once more.
\end{tcolorbox}
